\subsection{Key results}
\label{sec:discussion-key}
This national, two-phase-recruited sample of competitive Italian athletics athletes shows carbon-plated spikes established as the prevailing training equipment: a mean 35.7\% of preparation-period and 47.7\% of competition-period training-type frequency, with adoption above 80\% in both periods and a consistent pattern across sex and age-group strata. Adoption was not clearly associated with injury prevalence in either period. Among competition-period users, the primary nonlinear share model showed higher adjusted injury prevalence at higher share, but the matched linear comparison showed that this association was not unique to the relative scale: an increase from 37.3\% to 61.4\% share (24.1 percentage points) corresponded to a 7.9-pp risk difference (95\% CrI 1.4 to 14.4; $P_{\mathrm{dir}}=99.0\%$), while an increase from 2.83 to 6.19 units of the summed ordinal carbon-frequency score (3.36 units, not unique sessions) corresponded to 8.3 pp (95\% CrI 1.6 to 15.9; $P_{\mathrm{dir}}=99.2\%$). Both preparation estimates remained uncertain. The adapted OSTRC-H2 items added an athlete-centered burden estimate beyond injury counts: median severity was 75 of 100 and was highest among athletes reporting injuries in both periods. A further original descriptive finding reinforces this pattern: among the 248 lower-limb injuries with a recorded type, tendinopathy (28.6\%) and bone stress injury or stress fracture (26.2\%) were the two leading injury types, together exceeding muscle strain (22.2\%)--a distribution markedly different from general athletics injury surveillance, where muscle injury typically predominates \citep{edouard2020injury}. Cross-tabulating type against location refines this further: tendinopathy was disproportionately concentrated at the foot, lower leg, and ankle relative to the injury sample as a whole, a distal pattern that lines up with the Achilles- and forefoot-tendon loading flagged by the biomechanical literature on carbon-plate use (Section~\ref{sec:discussion-comparison}); bone stress injury, by contrast, was thigh-dominant like the sample overall, so its elevated frequency in this cohort is a genuine finding but is not, on the present data, specifically localized to the foot or ankle.

\subsection{Interpretation}
\label{sec:discussion-interpretation}
The distinction between adoption and amount of use remains central: the association tracked greater competition-period exposure among users, not simply adoption. The exposure-scale comparison, however, changes the mechanistic interpretation. Similar standardized estimates for share and absolute frequency mean that these data cannot identify whether relative allocation to carbon footwear or cumulative carbon-use frequency is the more relevant quantity. The two measures are mathematically related and were therefore fitted separately; adjustment for period-matched total training frequency makes each coefficient conditional on reported training volume but does not turn their descriptive comparison into causal separation. The result is compatible with a cumulative-load mechanism in which greater exposure to altered longitudinal bending stiffness changes ankle and plantarflexor loading \citep{gerlach2025aft,willwacher2016stiffness}, but it is also compatible with unmeasured competition-period intensity, footwear allocation, or reverse causation.

The primary HSGP curve, its linear sensitivity analysis, and the matched standardized models agree on the positive direction of the competition association, at a magnitude estimated in the 6--15 pp range across specifications. The empirical P75-versus-P25 estimate is a summary of one continuously estimated dose--response curve, not evidence for a threshold or for share as the clinically privileged scale. Likewise, the exact 25\%, 50\%, and 75\% values shown on the profile-specific curve are interpretive anchors on that same curve: no athlete groups were formed at those values, and the model did not estimate change-points at 25\%, 50\%, or 75\%. The profile-specific outputs add clinical interpretability by expressing fitted associations as absolute injury probabilities while holding measured characteristics constant; the same models can generate predictions for other specified profiles, and validating them as individual risk calculators is a task for future, externally validated work. The appropriately scoped conclusion is that greater reported competition-period carbon exposure among users was associated with greater adjusted injury prevalence on both examined scales--a population-level signal worth acting on through monitoring, even though its precise magnitude is not yet settled. Prospective measurement of unique carbon-shoe sessions, duration, intensity, and temporal order is needed to distinguish relative from cumulative exposure and to assess causality.

The OSTRC findings also show why injury burden should not be reduced to absence or complete cessation from sport. Time-loss-only systems can underestimate overuse problems that impair training, performance, or symptoms while athletes continue to participate \citep{clarsen2013ostrc,clarsen2020ostrc}. Direct athlete reporting therefore complements injury counts and time loss with consequences that matter to the athlete. In this study the adapted score is best understood as a sport-specific athlete-reported outcome measure, not as a generic quality-of-life instrument or a substitute for prospectively repeated OSTRC surveillance.

\subsection{Comparison with prior evidence}
\label{sec:discussion-comparison}
This study adds a population-level description to a literature that has so far relied on biomechanical and case-based work. A kinetic comparison of traditional and carbon-plated sprint spikes reported higher ankle joint contact force and peak plantarflexor force in the carbon-plated condition, with substantial inter-individual variability \citep{gerlach2025aft}, providing a mechanistic account well aligned with the dose-sensitive pattern reported here. Case-based literature on Achilles tendon load with carbon-plate use \citep{tenforde2023bonestress} points in the same direction as the tendinopathy pattern reported here: tendinopathy was both over-represented as an injury type relative to the typical athletics pattern \citep{edouard2020injury} and disproportionately located at the foot, lower leg, and ankle (Table~\ref{tab:injury-type-by-location}), consistent with that case-based signal at the population level for the first time. The same case-based literature has also flagged navicular and cuneiform bone stress with carbon-plate use \citep{tenforde2023bonestress}; this cohort's elevated overall rate of bone stress injury is compatible with that concern in aggregate, but the injuries were thigh-dominant rather than foot-dominant in location, and the survey did not resolve bone-specific site (for example, navicular versus tibial), so this cohort's data support the tendinopathy hypothesis more directly than the site-specific bone-stress-injury hypothesis. To the authors' knowledge, this is among the first studies to describe carbon-plate insole exposure and its association with injury prevalence at the population level in a national sample, extending this prior mechanistic and case-based work with cohort-level evidence; the tendinopathy and bone-stress-injury patterns observed here are hypothesis-generating findings that a dedicated, adequately powered type-stratified analysis should test directly.

\subsection{Strengths and considerations for interpretation}
\label{sec:discussion-limitations}
The two-phase recruitment, national age--sex poststratification, and Bayesian multilevel design were built specifically to support a defensible population-level estimate from survey data, and several features of the results reflect that: consistent direction across sensitivity specifications, well-behaved sampling diagnostics, and posterior predictive checks that closely reproduced observed injury counts. A number of features of the data are useful to keep in mind when applying these results. Because exposure and outcome were recalled together for the same season, this cross-sectional design describes an association rather than establishing temporal order between carbon-exposure increase and injury onset; a prospective design with objectively measured exposure is the natural next step. The survey's timing compounds this: administered near the end of the season, it asked athletes to recall the competition period from recent memory but the preparation period from further back, a differential recall window that could itself help explain why the preparation-period association came through weaker and less certain than the competition-period one (Section~\ref{sec:bias}). Both exposure scales were derived from ordinal weekly training-type-frequency items. In particular, the absolute score is not an exact count of unique sessions: responses were top-coded and training domains may overlap within a session. The scale comparison was added to clarify interpretation after the primary analysis and should be regarded as secondary. The competition-period non-user group, while smaller (44 of 352 athletes) than the user group, reflects the near-universal adoption of carbon-plate spikes in this population and widens the credible interval for the adoption contrast. The association magnitude varied with functional form and hyperprior width, so the primary HSGP estimate of 13.2 pp is evidence of a positive association of plausibly meaningful size, with a magnitude this study's data place somewhere in a 6--15 pp range. The adjustment set covered principal risk factors identified in prior running-injury literature \citep{taunton2003running}; coaching-driven equipment allocation, session intensity, and baseline plantarflexor strength were not measured. Self-report enabled national-scale surveillance \citep{bahr2020strobesiis}, but objective training logs are needed to validate the exposure measures. The adapted OSTRC-H2 items were completed by 94.6\% of injured athletes but were asked only once and only after injury was reported, so they cannot recover weekly prevalence, symptom duration, cumulative burden, or scores among athletes who did not classify themselves as injured; the period-stratified SMDs are correspondingly descriptive. The study protocol was prospectively reviewed and approved by the Comitato Etico of the University of Turin (Section~\ref{sec:registration}), providing a documented, ethics-verified record of pre-specification alongside this manuscript.

\subsection{Generalizability}
\label{sec:discussion-generalizability}
The two-phase recruitment and age--sex poststratification were designed specifically to support generalizability to the national competitive population; the modest Kish effective-sample-size retention after weighting (97\%) and the close agreement between weighted and unweighted estimates (Appendix Table~\ref{tab:weighted-vs-unweighted}) support good generalizability to Italian FIDAL-affiliated competitive athletes across the age categories studied. Extending these findings to recreational athletes, road-running populations, or national contexts with different equipment-access patterns is a promising direction for future work, particularly given the near-universal adoption of carbon-plate spikes observed in this competitive cohort.

\subsection{Conclusion}
\label{sec:discussion-conclusion}
In this national cohort of competitive Italian athletics athletes, carbon-plated spikes were the prevailing training equipment in both preparation and competition. Adoption was not clearly associated with injury prevalence. Among competition-period users, greater reported carbon exposure was associated with higher adjusted injury prevalence whether represented as share or absolute frequency score; preparation-period estimates remained uncertain, and the data did not identify one exposure scale as superior. These findings support prospective surveillance using objectively measured unique sessions, duration, intensity, and known temporal order. Until such evidence is available, athletes, coaches, and clinicians should interpret the observed association as a signal for monitoring rather than as a causal risk estimate or a validated threshold for limiting carbon-plate use.
